We then ran a 2$\times$2 study that crosses action conditioning (an action-conditional predictor versus an action-free ablation) with the collection method (branching versus no branching). Each collection method had \StudyCfgPretrainBudget{} pretraining transitions per seed, and every encoder was read out through a separate restore-free pool. The readout again recognised the informative control, at forced-commit accuracy \StudyForcedInfoMin{} in all seeds. None of the four learned encoders exceeded the untrained encoder: their accuracies ranged from \StudyForcedLearnedMin{} to \StudyForcedLearnedMax{}, so the pre-registered learning-benefit contrast failed (\StudyVerdict). Because every learned arm stayed at chance, the collection and conditioning contrasts sit at a floor and do not separate the two effects; we do not read them as evidence of no effect.